\reportheading{附录 A.\quad 工况登记}
第 18 组在 OP-1 和 OP-6 下均取 47.2 和 106.2~kW，依次对应文献~\cite{peng} 的实验编号 4、8、44、48。敏感性试验以这些基本工况为基础，分别调整网格、亚网格模型、辐射份额或壁面处理。丙烷火剖面为 FDS 6.10.1 在 90--120~s 的时间加权平均。丙烯火筛选在首次观测到越限或达到 120~s 时结束。
\subhead{附录 B.\quad 网格与门洞流量对照}
47.2 和 106.2~kW 的火源特征直径分别为 0.283 和 0.391~m。两者与单元最长边之比在粗、中、细网格上分别为 5.65/7.82、7.27/10.05 和 9.42/13.03。与网格对齐后的燃烧器面积分别为 0.06250、0.06125 和 0.06300~m$^2$。
\tabletitle{表 B1.\quad 门流诊断工况，90--120~s。}
\begin{reporttable}{L{.16}L{.22}L{.21}L{.21}L{.20}}
\rowcolor{tablegrey}门态&火源 (kW)&进气 (kg/s)&排气 (kg/s)&RMSE (K)\\
OP-1&47.2&0.260&0.261&166.6\\
OP-1&106.2&0.318&0.322&55.0\\
OP-6&47.2&0.094&0.096&61.4\\
OP-6&106.2&0.095&0.098&361.4\\
\end{reporttable}
\reportfigure{\begin{tikzpicture}
\begin{axis}[width=.487\linewidth,height=5.8cm,axis lines*=left,ymin=0,ymax=.65,
ytick={0,.15,.30,.45,.60},xmin=-1.2,xmax=3.6,
xlabel={向外速度 (m/s)},ylabel={高度 (m)},ymajorgrids,grid style={ink!13},
tick label style={font=\sffamily\footnotesize},label style={font=\sffamily\small},
legend style={at={(.03,.97)},anchor=north west,font=\sffamily\footnotesize,draw=none}]
\addplot[black,densely dashed,forget plot] coordinates {(0,0) (0,.63)};
\addplot[teal,mark=*,mark size=1.7pt] coordinates {(-0.6169,0.06) (-0.5525,0.14) (-0.4858,0.22) (-0.1781,0.3) (0.3791,0.38) (0.8191,0.46) (1.459,0.54) (1.8563,0.6)};
\addlegendentry{OP-1}
\addplot[rust,mark=triangle*,mark size=2.2pt] coordinates {(-0.8488,0.06) (-0.6996,0.14) (-0.5641,0.22) (0.1562,0.3) (0.6777,0.38) (1.6367,0.46) (3.164,0.54) (3.2764,0.6)};
\addlegendentry{OP-6}
\end{axis}\end{tikzpicture}%
\begin{tikzpicture}
\begin{axis}[width=.487\linewidth,height=5.8cm,axis lines*=left,ymin=0,ymax=.65,
ytick={0,.15,.30,.45,.60},xmin=0,xmax=950,
xlabel={温度 ($^\circ$C)},ylabel={高度 (m)},ymajorgrids,grid style={ink!13},
tick label style={font=\sffamily\footnotesize},label style={font=\sffamily\small},
legend style={at={(.03,.97)},anchor=north west,font=\sffamily\footnotesize,draw=none}]

\addplot[teal,mark=*,mark size=1.7pt] coordinates {(35.0545,0.06) (29.7306,0.14) (29.8055,0.22) (59.4344,0.3) (115.7244,0.38) (225.2291,0.46) (353.3391,0.54) (482.9451,0.6)};
\addlegendentry{OP-1}
\addplot[rust,mark=triangle*,mark size=2.2pt] coordinates {(51.0932,0.06) (38.5634,0.14) (42.8348,0.22) (247.3595,0.3) (467.4785,0.38) (702.2993,0.46) (838.2762,0.54) (890.311,0.6)};
\addlegendentry{OP-6}
\end{axis}\end{tikzpicture}}{图 B1.\quad 106.2~kW 绝热诊断工况中 R1 门洞八个高度处的速度与温度。各点为 90--120~s 平均值，负速度表示流入房间。}

\clearpage
\reportheading{附录 C.\quad 原始 OP-6 方案搜索}
\tabletitle{表 C1.\quad 不同垂壁深度和排烟量下的采样净高余量（mm）。}
\begin{reporttable}{L{.28}L{.18}L{.18}L{.18}L{.18}}
\rowcolor{tablegrey}垂壁深度&\multicolumn{4}{c}{总排烟量 (m$^3$/s)}\\
(m)&0.08&0.16&0.24&0.32\\
0.00&$-334$&$-331$&$-295$&$-289$\\
0.08&$-336$&$-331$&$-288$&$-294$\\
0.16&$-335$&$-331$&$-296$&$-285$\\
0.24&$-336$&$-337$&$-294$&$-302$\\
\end{reporttable}
\clearpage
\reportheading{附录 D.\quad 风量对照与数据覆盖}
\tabletitle{表 D1.\quad 原网格八个测点在 0--120~s 内的最低热分层界面高度（m）。目标值为 0.380~m。}
\begin{reporttable}{L{.23}L{.23}L{.27}L{.27}}
\rowcolor{tablegrey}门态&烟炱产率 (kg/kg)&0.96 m$^3$/s&1.28 m$^3$/s\\
OP-1&0.02&0.4574&\textbf{0.4425}\\
OP-1&0.05&0.4717&0.5093\\
OP-1&0.10&0.4410&0.5013\\
OP-6&0.02&\textbf{0.3957}&0.5827\\
OP-6&0.05&0.5262&0.4975\\
OP-6&0.10&0.4061&0.5503\\
\end{reporttable}
表 D1 中 1.28~m$^3$/s 的最低值已与原始设备记录和正常完成日志核对。0.96~m$^3$/s 的六组工况均有按测点归档的数据摘要，覆盖八个测点和 0--120~s。低风量比较采用上述全时段归档数据，部分本地原始文件对应早期中断的计算。在 1.28~m$^3$/s 下，OP-1 基准烟炱产率的细网格最低值（0.5879~m）也来自八测点归档摘要，本地原始记录止于 20~s。

\tabletitle{表 D2.\quad 改进布局、OP-6、烟炱产率 0.02、竖向 38 层网格。最低值取自记录时段的全部八个测点。}
\begin{reporttable}{L{.20}L{.22}L{.22}L{.36}}
\rowcolor{tablegrey}$Q$ (m$^3$/s)&输入功率 (W)&最低高度 (m)&FDS 记录\\
0.96&4.35&0.3717&98.0 s 越限\\
1.12&6.91&0.5522&完整，120 s\\
1.28&10.32&0.5429&完整，120 s\\
\end{reporttable}

合计输入功率按第 2.2 节的两条短支管、效率 0.60 计算。表 D2 中功率最低的可行风量为 1.12~m$^3$/s。表 D3 给出该风量六组完整计算各自的八测点最低值。
\tabletitle{表 D3.\quad 1.12~m$^3$/s 的细网格复核。净高阈值为 0.38~m。}
\begin{reporttable}{L{.16}L{.24}L{.14}L{.23}L{.23}}
\rowcolor{tablegrey}门态&烟炱产率 (kg/kg)&状态&记录时长 (s)&最低高度 (m)\\
\csname @@input\endcsname sections-zh/intermediate-coverage.tex
\end{reporttable}
\clearpage
\reportheading{附录 E.\quad 代理模型的越限事件评价}
本研究按 PIELM 方法~\cite{pielm} 建立固定隐层，采用 96 个固定的 tanh 单元，输出层以岭系数 0.1 作最小二乘拟合。参考区域模型按气层建立控制体收支的思路~\cite{cfast}，本研究假设横向均匀，用简化的上层体积平衡约束加权平均层厚。短时事件定义为单个测点连续低于 0.38~m 且持续不超过 2~s。预测与观测时段的交集至少占并集的一半时，按一一对应原则计为匹配。设备记录间隔约 0.5~s，事件的起止时刻由相邻记录线性插值得到。
\subhead{直接检验最低高度判定}
在较早的 20 个算例（竖向 19 层网格）上，拟合测点最低值修正项后，按完整设计留出的最低高度误差从 27.03 降至 8.26~cm。八个越限算例的漏判数从八降至六，十二个通过算例中有一次误报。扩展样本的比较见第 5 节。

对 OP-6 低风量细网格算例，计算结束前冻结的修正预测为 0.4710~m。98~s 时记录到 0.3717~m 后即停止，故预测至少高估 9.9~cm。对 OP-6、基准烟炱产率、1.12~m$^3$/s 的算例，60 和 90~s 时发出的预测均识别出 $x=4.2$~m 的控制测点，该测点的最低值出现在 49.50~s。

另一项提前两秒的预警试验使用计算过程中的观测。事后将两模型的误报数均调至 33 次时，局部高度模型与多测点温度模型在 25 个短时事件中检出了相同的四个，平均提前 0.73~s。
\newcommand{\stationtrace}[2]{\begin{tikzpicture}
\begin{axis}[width=.48\linewidth,height=2.47cm,xmin=0,xmax=120,ymin=0,ymax=.80,xtick={0,60,120},ytick={0,.38,.76},scaled ticks=false,title={$x=#2$ m},ymajorgrids]
\addplot[ink!75,densely dashed] coordinates {(0,.38) (120,.38)};
\addplot[teal,thick] table[col sep=comma,x=time,y=fds#1]{data/pinn-eight-stations.csv};
\addplot[rust,dashed,thick] table[col sep=comma,x=time,y=pinn#1]{data/pinn-eight-stations.csv};
\end{axis}\end{tikzpicture}}
\reportfigure{{\sffamily\footnotesize 烟层高度 (m)\hfill\textcolor{teal}{\rule[.4ex]{12pt}{1pt}} FDS\quad\textcolor{rust}{-- --} PINN\quad\textcolor{ink}{-- --} 0.38 m 阈值\par}
\begin{tabular}{@{}cc@{}}
\stationtrace{0}{0.5}&\stationtrace{1}{1.7}\\
\stationtrace{2}{2.4}&\stationtrace{3}{3}\\
\stationtrace{4}{3.6}&\stationtrace{5}{4.2}\\
\stationtrace{6}{4.9}&\stationtrace{7}{6.1}
\end{tabular}\par{\sffamily\footnotesize 时间 (s)}}{图 D1.\quad FDS 与 PINN 在八列测点处的热分层高度时序。远端测点的局部放大见图 D2。}

\clearpage
\reportheading{附录 F.\quad 敏感性测试与按布局留出验证}
\tabletitle{表 F1.\quad 106.2~kW 的 OP-6 板材模型与 文献~\cite{peng} 图 9 的对照。峰位以燃烧器中心为原点。}
\begin{reporttable}{L{.17}L{.21}L{.20}L{.23}L{.19}}
\rowcolor{tablegrey}时间窗 (s)&峰值温升 (K)&峰位 $x$ (m)&R1 侧偏差 (K)&RMSE (K)\\
90--105&858.2&$-0.40$&$-119.9$&107.8\\
105--120&901.3&$-0.30$&$-105.1$&95.0\\
90--120&876.0&$-0.30$&$-112.5$&100.7\\
\end{reporttable}
\tablenote{实测峰值为 890.8~K，位于 $-0.20$~m。偏差为 $+1.0$ 至 $+3.3$~m 的 24 个测点上模拟减实验的均值。RMSE 使用全部 51 个测点。正文统一采用 90--120~s 时间窗。}

\medskip
\tabletitle{表 F2.\quad 逐一将整类补风布局排除在训练集之外的最低高度预测。分类阈值为 0.38~m。}
\begin{reporttable}{L{.31}L{.23}L{.23}L{.23}}
\rowcolor{tablegrey}估计模型&MAE (cm)&越限漏判&通过误报\\
修正 PINN&9.43&14/14&0/19\\
岭回归&8.87&14/14&0/19\\
最近邻&9.59&14/14&0/19\\
\end{reporttable}
\tablenote{MAE 取 27 次完整计算。分类另纳入六次观测到越限后停止的计算，共 33 次判定。按单个设计留出的结果另见第 5 节。}

\medskip
\subhead{低功率 OP-1 温度诊断}
新增探头复现了基准热电偶的 192.4~K RMSE。用气体温度直接与 Peng 实验对照，RMSE 为 187.6~K。气体平均比热电偶测珠高 4.7~K，占 191.8~K 平均低估量的 2.5\%。在相邻两个 15~s 时间窗内，热电偶平均分别低估 191.1 和 192.5~K。这组同位置对照量化了气体与测珠的温差。羽流混合与门洞换气的贡献未在这组对照中单独测定。
\reportfigure{\begin{tikzpicture}
\begin{axis}[width=.95\linewidth,height=4.7cm,xmin=-340,xmax=340,
xtick={-300,-150,0,150,300},xlabel={距火源轴线距离 (cm)},
ylabel={温升 (K)},ymajorgrids,legend columns=3,
legend style={at={(.5,1.02)},anchor=south}]
\addplot[ink,only marks,mark=*,mark size=1.1pt] table[col sep=comma,x=figure_r_cm,y=exp_rise_K]{../../05_Calculations/op1_47_plume_diag_stations.csv};\addlegendentry{Peng 等}
\addplot[teal,thick] table[col sep=comma,x expr=\thisrow{model_x_cm}-330,y=gas_rise_K]{../../05_Calculations/op1_47_plume_diag_stations.csv};\addlegendentry{FDS 气体}
\addplot[rust,thick,dashed] table[col sep=comma,x expr=\thisrow{model_x_cm}-330,y=fds_rise_K]{../../05_Calculations/op1_47_plume_diag_stations.csv};\addlegendentry{FDS 热电偶}
\end{axis}\end{tikzpicture}}{图 F1.\quad 47.2~kW、OP-1。绝热粗网格上，在顶棚下 20~mm 处同步采集气体与 1~mm 热电偶温度。实测曲线取自文献~\cite{peng} 图 9。模拟温升为 90--120~s 时间加权平均。}

